\subsection{三点を放物線で結ぶ}

次に，区間の両端に加えて中点の値も使う。中点を$m$，半分の幅を$h$とし，区間を$[m-h,m+h]$と書く。$y=x-m$によって中心を$0$へ移して，三点を通る二次式を
\[
 q(y)=A+By+Cy^2
\]
とする。$q(0)=f(m)$より$A=f(m)$である。また，$q(h)$と$q(-h)$を足せば$B$の項が消えるので
\[
 C=\frac{f(m-h)-2f(m)+f(m+h)}{2h^2}.
\]
積分でも$By$は奇関数として消える。したがって
\[
\begin{aligned}
 \int_{-h}^hq(y)\dd y
 &=2hA+\frac{2h^3}{3}C\\
 &=\frac h3\{f(m-h)+4f(m)+f(m+h)\}.
\end{aligned}
\]
これがシンプソン公式である。中点の係数$4$は，二次式の積分を計算した結果として決まる。

この公式は三次式についても正確である。中心を移した三次式のうち，三次と一次の項は奇関数なので，積分も両端と中点の重み付き和も$0$になる。残った二次と定数の項については，すでに正確である。この「三次式まで誤差がない」という性質により，誤差は四階微分で評価できる。

